% 仙后座（综述）
% license CCBYSA3
% type Wiki

本文根据 CC-BY-SA 协议转载翻译自维基百科\href{https://en.wikipedia.org/wiki/Cassiopeia_(constellation)}{相关文章}。

仙后座（Cassiopeia，ⓘ）是位于北天的一個星座與星群，其名稱來自希臘神話中虛榮的王后卡西奧佩婭（Cassiopeia），她是安德洛墨達（Andromeda）的母親，曾誇耀自己擁有無與倫比的美貌。仙后座是公元 2 世紀希臘天文學家托勒密列出的 48 個星座之一，至今仍是現代 88 個星座之一。由於其由五顆亮星組成的獨特「W」形狀，仙后座極易辨認。

仙后座位於北天，在北緯 34° 以上的地區全年可見；在（亞）熱帶地區，9 月至 11 月初觀測條件最佳；而在南半球低緯度、低於南緯 25° 的熱帶地區，則可於特定季節在北方低空看到。

視星等為 2.2 的仙后座 α 星，即舍達爾（Schedar），是仙后座中最亮的恆星。該星座擁有一些已知最為明亮的恆星，包括黃色特超巨星仙后座 ρ 星（Rho Cassiopeiae）與 V509 Cassiopeiae，以及白色特超巨星仙后座 6 星。1572 年，第谷·布拉赫（Tycho Brahe）所觀測到的超新星就在仙后座中明亮爆發$^{[5]}$。仙后座 A 是一個超新星殘骸，在頻率高於 1 GHz 時，是全天最亮的銀河系外射電源。已有 14 個恆星系統被發現擁有系外行星，其中之一——HD 219134——被認為可能擁有六顆行星。一段富集的銀河帶穿過仙后座，其中包含多個疏散星團、年輕而明亮的銀河盤恆星以及星雲。IC 10 是一個不規則星系，是目前已知距離最近的星暴星系，也是本星系群中唯一的此類星系。
\subsection{神話}
\begin{figure}[ht]
\centering
\includegraphics[width=6cm]{./figures/99da126f8a82dd17.png}
\caption{《乌拉尼娅之镜》中描绘的坐在宝座上的仙后座（Cassiopeia）} \label{fig_Cape_1}
\end{figure}
该星座以埃塞俄比亚王后卡西奥佩娅（Cassiopeia）命名。卡西奥佩娅是埃塞俄比亚国王刻甫斯（Cepheus）的妻子$^{[5]}$，也是公主安德洛墨达（Andromeda）的母亲。刻甫斯与卡西奥佩娅以及安德洛墨达一同被安置在星空中，相互毗邻。她之所以被放置在天空中，作为惩罚，是因为她夸口说自己的女儿安德洛墨达比涅瑞伊得斯（Nereids，海中女神）更加美丽，或者另一种说法是，她自称比海中仙女本身更加美丽，从而激怒了波塞冬$^{[6]}$。她被迫坐在宝座上绕北天极旋转，一半时间必须紧紧抓住宝座以免跌落；而波塞冬则裁定安德洛墨达应被绑在岩石上，作为海怪刻托斯（Cetus）的祭品。随后，安德洛墨达被英雄珀尔修斯所拯救，并最终与其成婚$^{[7][8]}$。

在历史上，卡西奥佩娅星座的形象有多种不同的描绘方式。在波斯，她被苏菲（al-Sufi）描绘为一位女王，右手持有一根带新月的权杖，头戴王冠，身旁还有一只双峰骆驼；在法国，她则被描绘为坐在大理石宝座上，左手持棕榈叶，右手提着自己的长袍。这一形象出自奥古斯丁·罗耶（Augustin Royer）于1679年出版的星图集$^{[7]}$。

在中国天文学中，构成仙后座的恒星分布于三个区域之中：紫微垣（紫微垣，Zǐ Wēi Yuán）、北方玄武（北方玄武，Běi Fāng Xuán Wǔ）以及西方白虎（西方白虎，Xī Fāng Bái Hǔ）。

中国天文学家在现代所称的仙后座区域内辨认出多种星象组合。κ、η 与 μ 仙后座星组成了名为“王桥”的星官；当它们与 α 与 β 仙后座星共同观测时，则构成了名为“王良”的大车。御者手中的鞭子由 γ 仙后座星表示，该星有时被称为“次”，即汉语中“鞭”的意思$^{[7]}$。

在印度神话中，仙后座（Cassiopeia）与神话人物沙尔米什塔（Sharmishtha）相关联——她是伟大的阿修罗（Daitya）国王弗里什帕尔瓦（Vrishparva）的女儿，也是德瓦雅尼（Devayani，对应安德洛墨达）的朋友。

在威尔士神话中，Llys Dôn（字面意思为“多恩的宫廷”）是该星座的传统威尔士名称。多恩的至少三位子女同样具有天文学上的对应：Caer Gwydion（“圭迪恩的堡垒”）是银河的传统威尔士名称，而 Caer Arianrhod（“阿丽安罗德的堡垒”）则对应北冕座（Corona Borealis）$^{[9]}$。

在17世纪，仙后座中的恒星曾被用来描绘多位《圣经》人物，其中包括所罗门之母拔示巴（Bathsheba）、《旧约》中的女先知底波拉（Deborah），以及耶稣的追随者抹大拉的玛利亚（Mary Magdalene）$^{[7]}$。

在一些阿拉伯星图集中，仙后座的恒星还被描绘为一个名为“染色之手”（Tinted Hand）的形象。据不同说法，这只手要么象征一只被指甲花染成红色的女子之手，要么象征穆罕默德之女法蒂玛（Fatima）沾满鲜血的手。该“手”由 α、β、γ、δ、ε 与 η 仙后座星组成；其“手臂”则由 α、γ、δ、ε、η 与 ν 英仙座星组成$^{[7]}$。

另一个包含仙后座恒星的阿拉伯星座是“骆驼”。其头部由仙女座的 λ、κ、ι 与 φ 星构成；驼峰是 β 仙后座星；身体由其余仙后座恒星组成；而双腿则由英仙座与仙女座中的恒星构成$^{[7]}$。

在其他文化中，人们还从这一星群中看到了手或麋鹿鹿角的形象$^{[10]}$。例如在萨米人（Sámi）的传统中，仙后座的 “W” 形状被视为麋鹿的鹿角；西伯利亚的楚科奇人（Chukchi）同样将其中五颗主星视作五只驯鹿雄鹿$^{[7]}$。

马绍尔群岛的居民将仙后座视为一幅巨大鼠海豚星座的一部分：仙后座的主星构成其尾部，仙女座与三角座（Triangulum）构成其身体，而白羊座则构成其头部$^{[7]}$。在夏威夷，α、β 与 γ 仙后座星各自有名称：α 仙后座星被称为 Poloahilani，β 仙后座星被称为 Polula，γ 仙后座星被称为 Mulehu。普卡普卡（Pukapuka）群岛的人们则将仙后座视为一个独立的星座，称为 Na Taki-tolu-a-Mataliki $^{[11]}$。
\subsection{特征}
\begin{figure}[ht]
\centering
\includegraphics[width=6cm]{./figures/8d913297d824a2aa.png}
\caption{夜空中的仙后座} \label{fig_Cape_2}
\end{figure}
仙后座曾发生过一次超新星事件，即仙后座A（Cassiopeia A），SN~1572。

仙后座覆盖面积为 $598.4$ 平方度，占全天面积的 $1.451\%$，在 $88$ 个星座中按面积排名第 $25$ 位$^{[12]}$。它北侧和西侧与仙王座（Cepheus）相邻，南侧和西侧与仙女座（Andromeda）相邻，东南侧与英仙座（Perseus）相邻，东侧与鹿豹座（Camelopardalis）相邻，并在西侧与蝎虎座（Lacerta）共享一小段边界。

该星座的三字母官方缩写由国际天文学联合会（International Astronomical Union）于 1922 年采用，为 ``Cas''$^{[13]}$。官方星座边界由比利时天文学家欧仁·德尔波特（Eugène Delporte）于 $1930$ 年制定$^{[b]}$，以一个由 $30$ 条边组成的多边形加以界定。在赤道坐标系中，这些边界的赤经范围为 $00^{\mathrm h}\,27^{\mathrm m}\,03^{\mathrm s}$ 至 $23^{\mathrm h}\,41^{\mathrm m}\,06^{\mathrm s}$，赤纬范围为 $+77.69^\circ$ 至 $+46.68^\circ$,$^{[3]}$。由于其位于北天球，该星座对南纬 $12^\circ$ 以北的观测者而言均可见$^{[12][c]}$。在高纬北方天空中，它属于拱极星座（即在夜空中永不落下），在不列颠群岛、加拿大以及美国北部地区均可全年观测到$^{[15]}$。
\subsubsection{特征}
\textbf{恒星}
\begin{figure}[ht]
\centering
\includegraphics[width=6cm]{./figures/88f670490ee393b8.png}
\caption{从北半球观测时，肉眼所见的仙后座星座} \label{fig_Cape_3}
\end{figure}
德国制图学家约翰·拜耳（Johann Bayer）使用希腊字母 Alpha 至 Omega，随后又使用 A 和 B，来标记该星座中最显著的 26 颗恒星。后来发现 Upsilon 实际上由两颗恒星组成，并由约翰·弗兰斯蒂德（John Flamsteed）标记为 Upsilon$^{[1]}$ 和 Upsilon$^{[2]}$。B Cassiopeiae 实际上就是著名的超新星——第谷超新星（Tycho's Supernova）$^{[16]}$。在该星座边界之内，共有 157 颗视星等小于或等于 6.5 的恒星$^{[d][12]}$。

“W”星群  
仙后座中最亮的五颗恒星——Alpha、Beta、Gamma、Delta 与 Epsilon Cassiopeiae——构成了具有代表性的 W 形星群$^{[15]}$。这五颗恒星均可用肉眼清晰观测，其中三颗具有明显的变光性，第四颗则被怀疑为低振幅变星。在北半球的春季与夏季夜晚，当该星群位于北极星下方时，其形状呈现为 W；而在北半球冬季，或从南半球纬度观测时，该星群位于北极星“上方”（即更接近天顶），其 W 形看起来则是倒置的。

Alpha Cassiopeiae，传统名称为 Schedar（源自阿拉伯语 Al Sadr，意为“胸部”），常被误认为是一个由四颗恒星组成的系统，但实际上它是一颗主星，伴随三颗在物理上距离遥远的视伴星。主星占据主导地位，是一颗橙色巨星，视星等为 2.2，距离地球$228\pm2$光年$^{[18]}$。其光度约为太阳的 771 倍，在其约 1 至 2 亿年的生命周期中，核心氢燃料耗尽后发生膨胀并冷却；在此前的大部分时间里，它是一颗蓝白色 B 型主序星$^{[19]}$。一颗视星等 8.9 的黄色矮星伴星（B）与其相距甚远，而另外两颗伴星（C 与 D）距离更近，视星等分别为 13 和 14$^{[20]}$。

Beta Cassiopeiae，又称 Caph（意为“手”），是一颗白色恒星，视星等为 2.3，距离地球 $54.7\pm0.3$ 光年$^{[18]}$。其年龄约为 12 亿年，核心氢已耗尽，正开始偏离主序阶段并发生膨胀和冷却。其质量约为太阳的 1.9 倍，光度约为太阳的 21.3 倍。Caph 的自转速度约为其临界速度的 92\%，完成一次自转仅需 1.12 天，这使得恒星呈现为扁球体形状，其赤道半径比极半径大约 24\%$^{[21]}$。它是一颗盾牌座$\delta$ 型变星，振幅较小，周期约为 2.5 小时$^{[22]}$。

Gamma Cassiopeiae，现称为 Tiansi，是 Gamma Cassiopeiae 型变星的原型，这类变星由于高速自转，会将物质抛射形成一个可变的盘状结构。Gamma Cassiopeiae 的最暗视星等为 3.0，最亮可达 1.6，但通常接近 2.2，其亮度变化具有不可预测的减弱与增强。它是一个分光双星系统，轨道周期为 203.59 天，其伴星的计算质量约与太阳相当。然而至今尚未获得该伴星的直接观测证据，因此有人推测它可能是一颗白矮星或其他简并天体$^{[23]}$。该系统距离地球 $550\pm10$ 光年。

Delta Cassiopeiae，也称 Ruchbah 或 Rukbat，意为“膝盖”，可能是一颗 Algol 型食双星，最大亮度约为 2.7 等。有报道称其存在振幅小于 0.1 等、周期约为 $2$ 年 $1$ 个月的食现象$^{[24]}$，但这一结果从未得到确认。它距离地球$99.4\pm0.4$光年$^{[18]}$。Delta Cassiopeiae 曾被让·皮卡尔（Jean Picard）于 1669 年用于测定地球表面上一度弧长，从而推算地球半径$^{[25]}$。

Epsilon Cassiopeiae，又称 Segin，视星等为 3.3，距离地球$410\pm20$光年$^{[18]}$。它是一颗炽热的蓝白色恒星，光谱型为 B3 III，表面温度约为$15{,}680\,\mathrm{K}$。其质量约为太阳的 6.5 倍，直径约为太阳的 4.2 倍，属于 Be 星这一类——高速自转并抛射出环状或壳状物质的恒星$^{[26]}$。

\textbf{较暗恒星}
\begin{figure}[ht]
\centering
\includegraphics[width=6cm]{./figures/f7c9e38a40b34540.png}
\caption{仙后座 $\kappa$ 星及其弓形激波。斯皮策红外望远镜图像（NASA/JPL-Caltech）} \label{fig_Cape_4}
\end{figure}
仙后座中接下来最亮的七颗恒星也全部被确认或被怀疑为变星，其中包括拜耳未赋予希腊字母的 50 仙后座星，它是一颗振幅极小的疑似变星。$\zeta$仙后座星（Fulu$^{[27]}$）是一颗疑似缓慢脉动的 B 型恒星。$\kappa$仙后座星（Cexing）是一颗光谱型为 BC0.7Ia 的蓝色超巨星，其光度约为太阳的 $302{,}000$ 倍，直径为太阳的 33 倍$^{[28]}$。它是一颗高速逃逸星，相对于其邻近恒星以约每小时 2.5 百万英里（每秒 $1{,}100$ 千米）的速度运动$^{[29]}$。其磁场与粒子风在其前方约 4 光年处形成了一个可见的弓形激波，与弥散且通常不可见的星际气体与尘埃发生碰撞。该弓形激波的尺度极其巨大：长度约 12 光年，宽度约 1.8 光年$^{[30]}$。$\theta$仙后座星，名为 Marfak，是一颗疑似变星，其亮度变化小于十分之一个星等。$\iota$仙后座星是一套距地球 142 光年的三星系统，主星是一颗视星等 4.5 的白色恒星，同时也是$\alpha_2$猎犬座型变星；次星是一颗视星等 6.9 的黄色恒星；第三星的视星等为 8.4。主星与次星彼此接近，而主星与第三星则相距甚远。$\omicron$ 仙后座星也是一套三星系统，其主星同样是一颗$\gamma$仙后座型变星。

$\eta$仙后座星（Achird）是仙后座中距离地球最近的恒星，距离约 19.13 光年$^{[31]}$。它是一套双星系统，由一颗与太阳相似、视星等 3.45 的 G 型矮星，以及一颗视星等 7.51 的 K 型矮星组成$^{[32]}$。它在天球上的位置介于$\alpha$与$\gamma$仙后座星之间，是距离地球最近的 G 型／K 型恒星之一。

$\sigma$仙后座星是一套距地球 $1{,}500$ 光年的双星系统，其主星呈绿色、视星等 5.0，次星呈蓝色、视星等 7.3。$\psi$ 仙后座星是一套距地球 193 光年的三星系统，主星是一颗视星等 4.7 的橙色巨星，次星为一对彼此接近、合成视星等约为 $9.0$ 的恒星$^{[24]}$。

$\rho$仙后座星是一颗半规则脉动变星型黄色特超巨星，其光度约为 $500{,}000\,L_{\odot}$，是银河系中最明亮的恒星之一$^{[33]}$。其最暗视星等为 6.2，最亮视星等为 4.1，周期约为 320 天。其直径约为太阳的 450 倍，质量为太阳的 17 倍，初始质量约为太阳的 45 倍。$\rho$ 仙后座星距地球约 $10{,}000$ 光年。仙后座还包括 V509 仙后座星，这是第二颗极为罕见的黄色特超巨星，其光度约为太阳的 $400{,}000$ 倍、质量为太阳的 14 倍$^{[33]}$，以及更高温的白色特超巨星 6 仙后座星。此外，该星座还包含红色超巨星 PZ 仙后座星，它是已知体积最大的恒星之一，估计半径为 $(1{,}190\text{--}1{,}940)\,R_{\odot}$，同时也是一颗半规则变星$^{[34]}$。其光度约为太阳的 $240{,}000$ 至 $270{,}000$ 倍，距离地球约 $9{,}160$ 光年$^{[35]}$。

AO 仙后座星是一套双星系统，由一颗 O8 型主序星和一颗 O9.2 型亮巨星组成，其质量分别约为太阳的 20.30--57.75 倍和 14.8--31.73 倍$^{[36]}$。这两颗大质量恒星彼此距离极近，因而相互引力作用使其形状被拉伸成近似卵形$^{[37]}$。

第谷·布拉赫的新星曾在仙后座中可见，而恒星 Tycho G 被怀疑是为该次爆炸提供触发物质的供体星。
\subsubsection{深空天体}
\begin{figure}[ht]
\centering
\includegraphics[width=6cm]{./figures/b85fc1bca2832008.png}
\caption{行星状星云 IC 289 是一团被恒星核心遗迹向外推入太空的电离气体云。} \label{fig_Cape_5}
\end{figure}
银河系中一段十分丰富的区域贯穿仙后座，从英仙座方向延伸至天鹅座，包含大量疏散星团、年轻而明亮的银河盘恒星以及各类星云。

心脏星云（Heart Nebula）和灵魂星云（Soul Nebula）是彼此相邻的两处发射星云，距离地球约 $7{,}500$ 光年。

仙后座中包含两个梅西耶天体，即 M52（NGC~7654）和 M103（NGC~581），二者均为疏散星团。M52 曾被描述为一个“肾形”的星团，约包含 100 颗恒星，距离地球约 $4{,}600$ 光年$^{[38]}$。其最显著的成员是一颗位于星团边缘、呈橙色、视星等为 8.0 的恒星。M103 的成员数量远少于 M52，仅约 25 颗恒星，同时其距离也更远，约为 $8{,}000$ 至 $9{,}500$ 光年$^{[39]}$。其最显著的成员实际上是一对距离更近、视向叠加的双星系统，由一颗 $7$ 等主星和一颗 10 等伴星组成$^{[24]}$。
\begin{figure}[ht]
\centering
\includegraphics[width=6cm]{./figures/fcc2a9e6beeaf1d2.png}
\caption{仙后座中的四个星团：IC~166、NGC~654、NGC~663 和 NGC~659} \label{fig_Cape_6}
\end{figure}
仙后座中另外一些显著的疏散星团包括 NGC~457 和 NGC~663，二者均约包含 80 颗恒星。NGC~457 的结构较为疏松，其最亮成员为仙后座$\phi$星（Phi Cassiopeiae），这是一颗呈白色、视星等为 5.0 的超巨星。然而，$\phi$仙后座是否真正属于该疏散星团目前尚不确定$^{[40]}$。NGC~457 中的恒星呈链状分布，距离地球约 $10{,}000$ 光年。NGC~663 的距离则更近，约为 $8{,}200$ 光年，同时其视直径更大，约为$0.25^\circ$,$^{[24]}$。

仙后座中存在两个超新星遗迹。第一个编号为 3C~10，亦称第谷超新星遗迹，是被称为“第谷之星”的超新星爆发之后的残留物。该超新星于 1572 年由第谷·布拉赫观测到，如今在射电波段中表现为一个明亮的天体$^{[24]}$。在由仙后座五颗主要恒星构成的 “W” 形星群内部，位于仙后座 A（Cas~A）。它是一次大约 300 年前发生的超新星爆发的遗迹（以地球当前观测到的时间计；其实际距离约为 $10{,}000$ 光年）$^{[41]}$，并且是太阳系之外可观测到的最强射电源之一。该天体可能曾于 1680 年被约翰·弗拉姆斯蒂德记录为一颗暗弱恒星。它也是钱德拉 X 射线天文台在 1990 年代末期获取的首批图像的研究对象之一。由该恒星抛射出的物质外壳以每秒约 $4{,}000$ 千米（$2{,}500$ 英里）的速度膨胀，其平均温度约为 $30{,}000\,\mathrm{K}$,$^{[41]}$。
\begin{figure}[ht]
\centering
\includegraphics[width=6cm]{./figures/8b22206cd9e7fd56.png}
\caption{仙后座，显示了双星团 $\chi$ 和 $h$ 英仙座，以及星团 NGC~654、NGC~663、M~103、NGC~457、NGC~225、NGC~7788、NGC~7790、NGC~7789 和 M~52} \label{fig_Cape_7}
\end{figure}
NGC~457 是仙后座中的另一处疏散星团，也被称为 E.T. 星团、猫头鹰星团以及 Caldwell~13。该星团于 1787 年由威廉·赫歇尔发现。其总星等为 6.4，距离地球约 $10{,}000$ 光年，位于银河系的英仙臂之中。然而，其最显著的成员——双星 $\phi$~Cassiopeiae——要近得多，距离仅约 $1{,}000$ 至 $4{,}000$ 光年。NGC~457 较为富集，属于 Shapley~e 级、Trumpler~I~3~r 型星团，其中心较为集中，并与周围恒星场分离。该星团包含超过 100 颗恒星，其亮度差异很大。$^{[42]}$

仙后座中还包含本星系群的两个成员。NGC~185 是一颗 9.2 等的 E0 型椭圆星系，距离地球约 200 万光年。稍暗且更远的 NGC~147 是一颗 9.3 等的椭圆星系，同样为 E0 型，距离地球约 230 万光年。尽管它们在天空中并不位于仙女座星系之内，但这两个矮星系都在引力上束缚于体量大得多的仙女座星系。$^{[43]}$

IC~10 是一颗不规则星系，是目前已知距离最近的星暴星系，也是本星系群中唯一的此类星系。$^{[44]}$

仙后座还包含距离本星系群最近的星系群之一——IC~342/Maffei 星系群的一部分。Maffei~1 和 Maffei~2 位于心脏星云与灵魂星云以南不远处。由于处在回避带区域，这两者尽管距离均在 $1{,}000$ 万光年以内，却显得异常昏暗（其中 Maffei~2 的亮度低于大多数业余望远镜的观测范围）。$^{[45]}$
\subsubsection{流星雨}  
十二月 $\phi$~仙后座流星雨是一场近年来发现的、发生在十二月初的流星雨，其辐射点位于仙后座。$\phi$~仙后座流星速度极慢，进入大气层时的速度约为每秒 16.7 千米。该流星雨的母体被认为是比拉彗星，但其母体的具体身份曾在相当长一段时间内并不明确。$^{[46]}$
\subsection{名称衍生}  
USS~Cassiopeia（AK-75）是一艘以该星座命名的美国海军陨石坑级货船。

在《宝可梦 朱／紫》中，反派组织“天星队”（Team~Star）被划分为五个小队，其名称均源自仙后座中最亮的恒星：Segin 小队、Schedar 小队、Ruchbah 小队、Navi 小队以及 Caph 小队。该组织的领导者使用“Cassiopeia”这一化名。

在《二之国：白色圣灰的女王》中，倒数第二位主要反派、亦即“白色女巫”，其名字正是卡西欧佩亚女王（Queen~Cassiopeia）。

Cassiopeia 也是伦敦乐队 Bears~in~Trees 的一首歌曲名称。尽管歌词主要指涉古希腊女性卡西欧佩亚，但该专辑封面展示了仙后座星座。$^{[47]}$

Cassiopeia 还是《英雄联盟》中的一名英雄角色。她的美貌与虚荣心映射了希腊神话中的同名人物。

Casiopea 是一个成立于 1976 年的日本爵士融合乐队的名称。该名称由吉他手的母亲选定，唯一的理由是便于乐队成员记住。
\subsection{参见}  
\begin{itemize}
\item 仙后座的恒星形成区
\end{itemize}  
\subsection{参考文献}  
\subsubsection{说明性注释} \\
a.$\gamma$~Cas 为变星，并且在某些时期亮度会超过 $\alpha$。\\
b.Delporte 曾向国际天文学联合会提议将星座边界标准化，该组织同意了这一提议，并授予他主导角色。$^{[14]}$\\  
c.尽管在南纬 $12^\circ$ 至 $43^\circ$ 之间的观测者技术上可以看到星座的部分区域升起至地平线上方，但距离地平线仅数度以内的恒星在实际上几乎无法观测。$^{[12]}$\\  
d.视星等为 6.5 的天体属于在郊区—乡村过渡夜空条件下肉眼可见的最暗天体之一。$^{[17]}$
\subsubsection{引用}
\begin{enumerate}
\item 琼斯，丹尼尔（2003）[1917]。彼得·罗奇；詹姆斯·哈特曼；简·塞特（编）。英语发音词典。剑桥：剑桥大学出版社。ISBN 3-12-539683-2.
\item “仙后座”. 莱克西科英国英语词典. 牛津大学出版社。于2020-03-22从原文存档。
\item “仙后座，星座边界”.星座. 国际天文学联合会. 检索日期：2016年12月2日
\item 柯克帕特里克，J·戴维；马罗科，费德里科等（2024年4月）。“基于全天空20秒差距范围内约3600颗恒星和棕矮星普查的初始质量函数”. 天体物理学报增刊系列. 271 (2): 55. arXiv:2312.03639. 文献标识码:2024ApJS..271...55K. doi:10.3847/1538-4365/ad24e2.
\item 奇肖姆，休, 主编（1911）。“仙后座”. 大英百科全书。第5卷（第11版）。剑桥大学出版社。第460页。
\item 陈，P.K.（2007）。星座图册：夜空中的恒星与神话。天空。第82页。ISBN 9781931559386.
\item 斯塔尔 1988, 第14–18页
\item 陈 2007, 第82–83页
\item 斯奎尔，查尔斯（2013）。凯尔特神话与传说。信使公司。ISBN 978-0-486-12209-0.
\item 普塔克，罗伯特（1998）。天空故事：古代与现代。纽约：新科学出版社。第104页。
\item 马克斯恩，莫德·伍斯特（1941）。晨星初升：波利尼西亚天文学述评。耶鲁大学出版社。第281页。文献标识码:1941msra.book.....M.
\item 伊恩·里德帕斯.“星座：仙女座—印第安座”星图故事自费出版。检索日期：2016年12月2日
\item 拉塞尔，亨利·诺里斯(1922)。“星座的新国际符号”。大众天文学. 30: 469. 文献标识码:1922PA.....30..469R.
\item 里德帕斯，伊恩.“星座边界：现代星座轮廓的形成过程”星象故事. 自行出版. 取自2016年6月1日
\item 阿诺德，H.J.P；多赫蒂，保罗；摩尔，帕特里克（1999）。恒星摄影图集。佛罗里达州博卡拉顿：CRC出版社。第20页。ISBN 978-0-7503-0654-6.
\item 瓦格曼，莫顿（2003）。失落的星辰：约翰内斯·拜耳、尼古拉斯·路易·德·拉凯耶、约翰·弗莱姆斯特德及其他多位天文学家星表中的失踪、遗失与疑难恒星。弗吉尼亚州布莱克斯堡：麦当劳与伍德沃德出版公司。第91–92. ISBN 978-0-939923-78-6.
\item 博特尔，约翰·E.（2001年2月）。“博特尔暗空等级”。《天空与望远镜》。天空出版公司。检索日期：2015年6月6日
\item 范·勒文，F.（2007）。“新依巴谷归算的验证”。天文学与天体物理学. 474 (2): 653–64. arXiv:0708.1752. 文献编码:2007A&A...474..653V. doi:10.1051/0004-6361:20078357. S2CID 18759600.
\item 詹姆斯·B·（吉姆）凯勒教授。“舍达尔（仙后座α星）”。伊利诺伊大学。于2010年3月27日从原始网页存档。检索日期：2010年2月22日。
\item 梅森，布莱恩·D.；怀科夫，加里·L.；哈特科普夫，威廉·I.；道格拉斯，杰弗里·G.；沃利，查尔斯·E.（2001）。“2001年美国海军天文台双星光盘。I. 华盛顿双星目录”. 天文学杂志. 122 (6): 3466. 文献编码:2001AJ....122.3466M. DOI:10.1086/323920.
\item 谢，X.；莫尼耶，J. D.；赵，M.；佩德雷蒂，E.；特罗，N.；梅朗，A.；滕·布鲁梅拉尔，T.；麦卡利斯特，H.；里奇威，S. T.（2011）。“更冷与更热：对仙后座β星和狮子座α星的干涉成像”。天体物理学杂志. 732 (2): 68. arXiv:1105.0740. Bibcode:2011ApJ...732...68C. doi:10.1088/0004-637X/732/2/68. S2CID 14330106.
\item 卡勒，詹姆斯·B.（吉姆）。“卡夫”。恒星。伊利诺伊大学。检索日期：2016年12月5日
\item 哈拉梅茨，P.；哈布达，P.；施特弗尔，S.；哈德拉瓦，P.；科尔恰科娃，D.；考布斯基，P.；克尔蒂奇卡，J.；库巴特，J.；什科达，P.；什莱赫塔，M.；沃尔夫，M.（2000）。“Be型恒星的性质与本质。第二十部分：仙后座γ双星性质及轨道要素”。天文学与天体物理学. 364: L85 – L88. arXiv:astro-ph/0011516. Bibcode:2000A&A...364L..85H.
\item 里德帕斯与蒂里昂 2001, 第106–108页。
\item 皮卡尔，J.；沃勒，R.（1688）。地球的测量：巴黎皇家科学院多位成员为此目的所作若干观测的记述。[让·皮卡尔《地球的测量》译本] 由理查德·沃勒从法文翻译。[动物自然史回忆录。第2部分。] R. 罗伯茨；由T. 巴塞特销售。检索日期：2025年1月21日。
\item 卡坦扎罗，G.（2013）。“北天Be型恒星样本中Hα与Hβ的光谱图集”。天文学与天体物理学. 550(A79): 18.arXiv:1212.6608. 文献标识码:2013A&A...550A..79C. doi:10.1051/0004-6361/201220357.
\item “命名恒星”。IAU.org。检索日期：2018年7月30日。
\item 塞尔，S. C.；普林贾，R. K.；马萨，D.；赖恩斯，R.（2008）。“银河系早期B型超巨星光学与紫外光谱的定量研究。I. 基本参数”。天文学与天体物理学. 481 (3): 777. arXiv:0801.4289. Bibcode:2008A&A...481..777S. doi:10.1051/0004-6361:20077125. S2CID 1552752.
\item 克拉文，惠特尼（2014年2月21日）。“仙后座卡帕星的弓形激波：一颗巨大而炽热的超巨星”。Phys.org。检索日期：2016年12月6日。
\item 佩里，C. S.；贝纳利亚，P.；布鲁克斯，D. P.；史蒂文斯，I. R.；伊塞基亚，N. L.（2012）。“E-BOSS：一项广泛的恒星弓激波巡天。I. 方法与首个目录”。天文学与天体物理学. 538: A108。arXiv:1109.3689. 文献标识码:2012A&A...538A.108P. DOI:10.1051/0004-6361/201118116. S2CID 62840857.
\item 柯克帕特里克，J·戴维；马罗科，费德里科；盖利诺，克里斯托弗·R；拉古，亚杜克里什纳；法赫蒂，杰奎琳·K；巴尔达莱斯·加格留菲，丹妮拉·C；舒尔，史蒂文·D；阿普斯，凯文；施耐德，亚当·C；迈斯纳，亚伦·M；库赫纳，马克·J；卡塞尔登，丹；斯马特，R。L. Casewell, S. L. Raddi, Roberto (2024-04-01）。“基于全天空20秒差距范围内约3600颗恒星和棕矮星普查的初始质量函数”。天体物理学报增刊系列271（2）：55。arXiv：2312.03639文献标识码2024ApJS..271...55Kdoi10.3847/1538-4365/ad24e2ISSN0067-0049
\item 费尔南德斯，J.；等（1998），“邻近目视双星的基本恒星参数：仙后座η、牧夫座XI、蛇夫座70和飞马座85”，天文学与天体物理学, 338: 455–464, 文献标识码:1998A&A...338..455F
\item 斯托瑟斯，理查德·B.（2012）。“黄色超巨星呈现长次级周期？”天体物理学报. 751 (2): 151. 文献标识码:2012ApJ...751..151S. DOI:10.1088/0004-637X/751/2/151. S2CID 121048201.
\item 勒维斯克，艾米丽·M.；马西，菲利普；奥尔森，K. A. G.；普莱兹，贝尔特兰；若塞林，埃里克；迈德，安德烈；梅内，乔治（2005年8月）。“银河系红色超巨星的有效温度标度：偏冷，但并未如我们所想的那么冷。”天体物理学杂志. 628 (2): 973–985. arXiv:astro-ph/0504337. Bibcode:2005ApJ...628..973L. doi:10.1086/430901. S2CID 15109583.
\item 久野健；浅木洋；今井宏；大山智（2013）。“利用甚长基线干涉仪H2O脉泽天体测量法测定红超巨星Pz Cas的距离与自行”。天体物理学报. 774 (2): 107. arXiv:1308.3580. Bibcode:2013ApJ...774..107K. doi:10.1088/0004-637X/774/2/107. S2CID 118867155.
\item 霍勒，M.M.；诺伊豪瑟，R.；舒茨，B.F.（2010）。“O型和B型恒星以及红超巨星的质量与光度”。天文新闻. 331 (4): 349. arXiv:1003.2335. 文献编码:2010AN....331..349H. doi:10.1002/asna.200911355. S2CID 111387483.
\item 天文学与宇宙论。剑桥大学出版社档案馆。1928年。第125页起。GGKEY：KFJRG3PWW14。
\item 吴振宇；等（2009年11月），“银河系中疏散星团的轨道”，皇家天文学会月报, 399 (4): 2146–2164, arXiv:0909.3737, Bibcode:2009MNRAS.399.2146W, doi:10.1111/j.1365-2966.2009.15416.x, S2CID 6066790.
\item 桑纳，J.；格费特，M.；布伦岑多夫，J.；施莫尔，J.（1999）。“疏散星团的光度学与运动学研究。I. NGC 581（M 103）”。天文学与天体物理学. 349: 448–456. arXiv:astro-ph/9908059. 文献标识码:1999A&A...349..448S.
\item 张晓波；罗春奎；傅景楠（2012）。“年轻疏散星团NGC 457中的B型变星”。天文学杂志. 144 (3): 86. 文献标识码:2012AJ....144...86Z. DOI:10.1088/0004-6256/144/3/86. S2CID 250804951.
\item 威尔金斯，杰米；邓恩，罗伯特（2006）。300个天体：宇宙的视觉参考（第1版）。纽约布法罗：萤火虫图书。ISBN 978-1-55407-175-3.
\item 利维 2005, 第92–93页。
\item 利维 2005, 第180–181页。
\item 尼德弗，大卫·L.；阿什利，特里莎；斯莱特，科林·T.；奥特，于尔根；约翰逊，梅根；贝尔，埃里克·F.；斯塔尼米罗维奇，斯涅扎娜；普特曼，玛丽；马耶夫斯基，史蒂文·R.；辛普森，卡罗琳·E.；尤特，艾娃；奥斯特鲁，汤姆·A.；伯顿，W·巴特勒（2013）。“近邻星爆矮不规则星系IC 10中存在相互作用的证据”。天体物理学杂志快报. 779(2): L15.arXiv:1310.7573. Bibcode:2013ApJ...779L..15N. doi:10.1088/2041-8205/779/2/L15. S2CID 119238691.
\item I. D. 卡拉琴采夫（2005）。“本星系群及其他邻近星系群”。天文期刊. 129 (1): 178–188. arXiv:astro-ph/0410065. 文献标识码:2005AJ....129..178K. DOI:10.1086/426368. S2CID 119385141.
\item https://www.amsmeteors.org/2024/12/meteor-activity-outlook-for-december-7-13-2024/
\item “仙后座”. Bandcamp. 检索日期2023-09-01.
\end{enumerate}
\subsubsection{通用及引用来源}
\begin{itemize}
\item 克劳斯，O；里克，GH；比尔克曼，SM；勒弗洛克，E；戈登，KD；江上，E；比金，J；休斯，JP；扬，ET；欣茨，JL；昆茨，SP；海因斯，DC（2005）。“仙后座A超新星遗迹附近的红外回声”。科学. 308 (5728): 1604–6. arXiv:astro-ph/0506186. Bibcode:2005Sci...308.1604K. doi:10.1126/science.1112035. PMID 15947181. S2CID 21908980.
\item 莱维，大卫·H.（2005）,深空天体, 普罗米修斯图书,ISBN 1-59102-361-0
\item 里德帕斯，伊恩；蒂里翁，威尔（2001）,恒星与行星指南, 普林斯顿大学出版社,ISBN 0-691-08913-2
\item 里德帕斯，伊恩；蒂里昂，威尔（2007）。恒星与行星指南。伦敦：柯林斯。ISBN 978-0-00-725120-9.
\item 斯塔尔，尤利乌斯·D·W.（1988）。天空中的新图案：恒星的神话与传说。麦当劳与伍德沃德出版公司。ISBN 978-0-939923-04-5.
\end{itemize}
\subsection{外部链接}
\begin{itemize}
\item 仙后座 (分类)
\item 深空星座摄影指南：仙后座
\item 可点击的仙后座
\item 星象故事——仙后座
\item <bold>瓦尔堡研究所图像数据库（约160幅中世纪及早期现代的卡西俄佩亚图像）</bold>
\end{itemize}